\section*{Bab \ref{ch:g3:addsub} --- Penjumlahan dan Pengurangan}

\begin{solution}{exo:g3:addsub:1}
$346 + 253 = 599$; \quad
$478 + 265 = 743$; \quad
$1\,536 + 2\,876 = 4\,412$.
\end{solution}

\begin{solution}{exo:g3:addsub:2}
$587 - 234 = 353$ (periksa: $353 + 234 = 587$).

$805 - 348 = 457$ (periksa: $457 + 348 = 805$).

$1\,000 - 463 = 537$ (periksa: $537 + 463 = 1\,000$).
\end{solution}

\begin{solution}{exo:g3:addsub:3}
$38 + 7 = 38 + 2 + 5 = 45$; \quad
$56 + 9 = 56 + 4 + 5 = 65$; \quad
$45 + 28 = 45 + 30 - 2 = 73$.
\end{solution}

\begin{solution}{exo:g3:addsub:4}
$340 + 220 = 560$; \quad $675 + 210 = 885$; \quad
$512 - 300 = 212$; \quad $840 - 220 = 620$.
\end{solution}

\begin{solution}{exo:g3:addsub:5}
Dari $47$ ke $50$: $3$; dari $50$ ke $63$: $13$; seluruhnya
$3 + 13 = 16$. Cara bersusun membenarkannya: $63 - 47 = 16$.
\end{solution}

\begin{solution}{exo:g3:addsub:6}
$492 + 218$: taksiran $500 + 200 = 700$; tepatnya $710$.

$913 - 388$: taksiran $900 - 400 = 500$; tepatnya $525$.
\end{solution}

\begin{solution}{exo:g3:addsub:7}
$348 + 296 = 644$. Sekolah itu punya $644$ murid.
\end{solution}

\begin{solution}{exo:g3:addsub:8}
$224 - 178 = 46$. Nora masih punya $46$ halaman untuk dibaca.
\end{solution}

\begin{solution}{exo:g3:addsub:9}
Satuan: $7 + \text{?} = 2$ tidak mungkin, jadi $7 + 5 = 12$: lubang
satuannya $5$, simpan $1$. Puluhan: $\text{?} + 2 + 1 = 6$: lubang
puluhannya $3$ (tanpa simpanan). Ratusan: $3 + \text{?} = 8$: lubang
ratusannya $5$. Jadi $337 + 525 = 862$, dan setiap lubang hanya
punya satu angka yang mungkin --- jadi hanya ada satu cara.
\end{solution}

\begin{solution}{exo:g3:addsub:10}
Terjual: $145 + 89 = 234$ croissant. Dipanggang: yang terjual
ditambah yang tidak terjual, $234 + 17 = 251$ croissant.
\end{solution}

\begin{solution}{exo:g3:addsub:11}
Dia benar: $500 - 200 = 300$, ditambah $1$ menjadi $301$; secara
bersusun, $500 - 199 = 301$ juga. Mengambil $200$ berarti membuang
\emph{satu terlalu banyak}, jadi satu harus dikembalikan ---
mengurangkan tetangga yang ramah lalu membetulkannya sering lebih
cepat daripada meminjam.

\end{solution}
